\chapter{Meniscus Repair and Meniscectomy}

\section*{Introduction \& Clinical Indications}
Meniscus repair and meniscectomy are arthroscopic surgical procedures aimed at addressing symptomatic tears of the meniscal cartilage within the knee joint. The menisci, which are C-shaped fibrocartilaginous structures, play a critical role in load distribution, shock absorption, and joint stability. When a meniscus is torn, patients often experience significant pain, swelling, and mechanical symptoms such as locking or catching, which can severely impair mobility and quality of life. The surgical interventions focus on either repairing the torn meniscus to restore its anatomical integrity or performing a partial meniscectomy to remove the damaged portion of the meniscus. The decision between repair and meniscectomy is influenced by factors such as the tear's location, size, pattern, chronicity, and the patient's age and activity level. Preserving functional meniscal tissue is essential to safeguard long-term knee health and reduce the risk of premature osteoarthritis.

\begin{itemize}
  \item Arthroscopic meniscal surgery
  \item Knee arthroscopy with meniscal repair
  \item Knee arthroscopy with partial meniscectomy
  \item Meniscal debridement
  \item Arthroscopic knee surgery for meniscal tear
  \item Meniscal fixation
\end{itemize}

The primary surgical principle in meniscal intervention is to alleviate symptoms caused by a torn meniscus while preserving as much functional meniscal tissue as possible. The menisci are vital for normal knee biomechanics, contributing to load transmission, shock absorption, and joint stability. A torn meniscus disrupts these functions, leading to pain and mechanical symptoms, and if untreated, may accelerate degenerative changes in the knee joint.

For **meniscal repair**, the goal is to re-approximate a repairable tear and preserve function. Vascular zone, tear pattern, tissue quality, chronicity, stability, and associated injuries all affect healing; a red-zone location alone does not guarantee repairability.

When tissue is irreparable and symptoms are attributable to an unstable fragment, a **partial meniscectomy** may remove only the minimum tissue needed to create a stable rim. Degenerative tears often respond best to exercise-based and nonoperative care, and a tear in the avascular zone is not by itself an indication for surgery.

The history of meniscal surgery has evolved significantly from radical excision to preservation techniques. In the late 19th and early 20th centuries, open surgical techniques often led to total meniscectomy, providing immediate symptom relief but predisposing patients to premature knee osteoarthritis due to the loss of meniscal function.

A pivotal shift occurred in the 1960s and 1970s with the work of **Dr. Robert W. Jackson** and **Dr. Lanny L. Johnson**, who advanced arthroscopy, allowing for minimally invasive visualization and treatment of knee injuries. Initially, arthroscopic meniscectomy became the standard treatment for meniscal tears. However, as the long-term consequences of meniscal loss became evident, there was a growing emphasis on preservation techniques.

In the 1980s, **Dr. Richard O'Connor** and others advanced arthroscopic meniscal repair techniques, demonstrating the feasibility and benefits of suturing torn menisci. This evolution continues today, with ongoing refinements in repair techniques and the exploration of biological augmentation strategies, such as platelet-rich plasma and stem cells, aimed at improving healing rates and functional outcomes.

Meniscus repair or meniscectomy is indicated for patients with symptomatic meniscal tears that have not responded to conservative management. Specific clinical indications include:

\begin{itemize}
  \item Persistent, localized knee pain along the joint line, exacerbated by activity or twisting motions.
  \item Mechanical symptoms such as recurrent locking, catching, or clicking within the knee joint.
  \item A sensation of the knee "giving way" or instability during weight-bearing activities.
  \item Recurrent knee effusion that does not resolve with conservative measures.
  \item MRI confirmation of a meniscal tear, particularly in the vascularized "red zone" or large, displaced tears.
  \item Failure of a comprehensive trial of non-operative treatment lasting 6-12 weeks.
  \item Acute, traumatic meniscal tears in younger individuals, where preservation is critical.
  \item Meniscal tears associated with significant knee injuries, such as ACL ruptures.
  \item Specific tear patterns like meniscal root tears that compromise stability.
\end{itemize}

Surgery is selected for specific acute displaced or repairable tears, persistent mechanical symptoms, or associated injuries after clinical and imaging assessment. Many degenerative tears are initially managed nonoperatively, and early surgery is not automatically required.

However, meniscal surgery can also serve as a **secondary or salvage procedure** in certain contexts. If an initial repair fails or if symptoms recur after a partial meniscectomy, revision surgery may be considered. In cases of chronic, degenerative meniscal tears, conservative treatment is often the first approach, with surgery becoming an option only if symptoms persist despite maximal non-operative efforts.

\section*{Relevant Surgical Anatomy}
The knee joint houses two menisci: the medial and lateral menisci, both of which are crucial for knee function. The **medial meniscus** is larger, C-shaped, and less mobile, firmly attached to the joint capsule and the medial collateral ligament (MCL). This attachment makes it more prone to injury. In contrast, the **lateral meniscus** is smaller, O-shaped, and more mobile, with a less extensive capsular attachment and a posterior connection to the popliteus tendon, allowing for greater movement during knee flexion and extension.

The menisci are divided into three vascular zones, which are critical for surgical decision-making:

\begin{itemize}
  \item **Red-Red Zone:** The outer one-third, well-vascularized by branches from the medial and lateral genicular arteries. Tears in this zone have the highest potential for healing.
  \item **Red-White Zone:** The middle one-third, with a mixed blood supply. Tears here have moderate healing potential.
  \item **White-White Zone:** The inner one-third, largely avascular. Tears in this zone have very limited healing potential and are typically managed with meniscectomy.
\end{itemize}

Nerve supply to the menisci is primarily from branches of the tibial, obturator, and common peroneal nerves, contributing to proprioception and pain sensation. Key ligamentous attachments include the coronary ligaments anchoring the menisci to the tibial plateau and the transverse meniscal ligament connecting the anterior horns. Surgical landmarks for arthroscopic portal placement include the patella, patellar tendon, tibial tubercle, and joint line, facilitating precise instrument access and visualization.

\section*{Preoperative Workup \& Preparation}
A thorough preoperative workup is essential to optimize surgical outcomes and minimize complications:

\begin{itemize}
  \item **Comprehensive Medical Evaluation:** A detailed medical history and physical examination assess overall health and identify comorbidities.

  \item **Diagnostic Imaging:** MRI confirms the diagnosis and characterizes the meniscal tear, while weight-bearing X-rays assess for osteoarthritis.

  \item **Laboratory Tests:** Routine blood tests, including CBC and coagulation studies, evaluate anemia and bleeding risks.

  \item **Anesthesia Consultation:** Patients are evaluated by the anesthesia team to discuss options and assess readiness.

  \item **Medication Review and Adjustment:** Current medications are reviewed, and blood thinners are typically discontinued prior to surgery.

  \item **NPO Status:** Patients are instructed to refrain from eating or drinking for a specified period before surgery.

  \item **Informed Consent:** The surgeon explains the procedure, risks, and benefits, obtaining formal consent.

  \item **Preoperative Instructions:** Patients receive instructions on antiseptic showering, transportation, and postoperative care.

  \item **Prophylactic Antibiotics:** Intravenous antibiotics are administered before incision to reduce infection risk.
\end{itemize}

\textbf{Factors requiring optimization or a different treatment plan:}
\begin{itemize}
  \item Active, uncontrolled infection within the knee joint.
  \item Severe osteoarthritis often explains pain better than the meniscus and may make isolated meniscal surgery ineffective; it is not a universal contraindication when another indication exists.
  \item Uncorrectable coagulopathy or severe bleeding disorder.
  \item Medically unstable patients unable to tolerate surgery.
\end{itemize}

**Relative Contraindications \& Precautions:**
\begin{itemize}
  \item Degenerative meniscal tears in older patients with minimal symptoms.
  \item Poor meniscal tissue quality that may not support repair.
  \item Large, complex tears in the avascular "white zone."
  \item Documented patient non-compliance with rehabilitation protocols.
  \item Concomitant ligamentous instability that may affect repair success.
  \item History of Complex Regional Pain Syndrome (CRPS).
  \item Significant peripheral vascular disease affecting healing.
  \item Morbid obesity increasing surgical risks.
\end{itemize}

\section*{Surgical Technique \& Procedure}
Meniscal surgery is primarily performed arthroscopically, allowing for minimally invasive access to the knee joint. The procedure begins with the administration of anesthesia, which may be general or regional. The patient is positioned supine, and a tourniquet is applied to achieve a bloodless field.

Small incisions, or portals, are made around the knee joint. An arthroscope is inserted through one portal, providing visualization of the joint interior. The knee is irrigated with sterile saline to maintain a clear field. Specialized instruments are introduced through other portals to perform the necessary maneuvers.

Key operative steps typically include:

\begin{enumerate}
  \item **Diagnostic Arthroscopy:** A thorough inspection of the knee joint is performed to confirm the meniscal tear and assess for concomitant injuries.
  
  \item **Tear Characterization and Preparation:** The meniscal tear is probed and assessed for type, size, location, and stability. For repair, the torn edges may be debrided to stimulate healing.

  \item **Meniscal Repair (if indicated):**
      \begin{itemize}
        \item The torn meniscal fragments are reduced and stabilized using specialized sutures or bioabsorbable fixation devices.
        \item Techniques include **Inside-out**, **Outside-in**, or **All-inside** repair.
      \end{itemize}

  \item **Partial Meniscectomy (if indicated):**
      \begin{itemize}
        \item Unstable portions of the torn meniscus are resected using arthroscopic instruments.
        \item The remaining meniscal rim is contoured to create a stable edge.
      \end{itemize}

  \item **Joint Lavage and Final Inspection:** The knee is irrigated to remove debris, and a final inspection confirms the stability of the repair or the smooth contour of the meniscectomy site.

  \item **Closure:** The arthroscopic portals are closed with sutures or adhesive strips, and sterile dressings are applied.
\end{enumerate}

\section*{Complications \& Risks}
Meniscal surgery carries potential risks, categorized into general surgical risks and procedure-specific risks.

**General Surgical Risks:**
\begin{itemize}
  \item **Anesthesia Complications:** Adverse reactions to anesthetic agents.
  \item **Bleeding:** Excessive intraoperative or postoperative bleeding.
  \item **Infection:** Surgical site infection requiring treatment.
  \item **DVT and PE:** Formation of blood clots leading to serious complications.
  \item **Scarring:** Formation of hypertrophic or keloid scars.
\end{itemize}

**Procedure-Specific Risks:**
\begin{itemize}
  \item **Nerve or Vascular Injury:** Damage to peripheral nerves or blood vessels.
  \item **Joint Stiffness (Arthrofibrosis):** Excessive scar tissue formation.
  \item **Persistent Pain:** Chronic knee pain despite successful surgery.
  \item **Failed Meniscal Repair (Re-tear):** The repaired meniscus may fail to heal.
  \item **Progressive Osteoarthritis:** Altered biomechanics may accelerate osteoarthritis.
  \item **Persistent or Recurrent Effusion/Swelling:** Chronic swelling in the knee joint.
  \item **Hardware Complications:** Issues related to fixation devices.
  \item **Complex Regional Pain Syndrome (CRPS):** Severe chronic pain condition.
  \item **Incomplete Symptom Resolution:** Surgery may not fully resolve symptoms.
\end{itemize}

\section*{Postoperative Care \& Recovery}
Postoperative recovery varies based on the procedure performed. Following meniscal surgery, patients are monitored in the Post-Anesthesia Care Unit (PACU) for pain management and initial recovery. Most patients are discharged on the same day.

For **partial meniscectomy**, recovery is generally quicker. Patients may bear weight as tolerated and initiate gentle range-of-motion exercises within days. A structured physical therapy program typically begins within the first week, with a gradual return to activities expected within 4-6 weeks.

In contrast, **meniscal repair** requires a more protected recovery. Weight-bearing is often restricted for 4-6 weeks, and physical therapy progresses gradually to protect the repair. Full activity, especially high-impact sports, may take 4-6 months or longer, emphasizing the importance of adherence to rehabilitation protocols.

During arthroscopic surgery, the surgeon documents intraoperative findings, including the location, pattern, size, and stability of the meniscal tear. The vascularity of the torn tissue is assessed, and the quality of the meniscal tissue is noted. Concomitant injuries to ligaments or cartilage are also documented.

For **partial meniscectomy**, resected tissue is sent for pathology, confirming the presence of fibrocartilaginous tissue and any degenerative changes. In cases of **meniscal repair**, the focus is on the successful apposition of torn fragments, with no tissue sent for pathology.

A successful postoperative course is characterized by a gradual resolution of symptoms and restoration of knee function. Patients may experience some pain and swelling initially, which should improve with conservative measures. For partial meniscectomy, pain typically subsides within 1-2 weeks, with a return to normal activities expected within a month.

For meniscal repair, the initial phase is more protected, with strict adherence to weight-bearing and range-of-motion restrictions. Over time, patients should experience increased strength, flexibility, and stability, with the ultimate goal of returning to pre-injury activity levels without significant limitations.

Postoperative care is crucial for monitoring recovery and addressing complications. Follow-up appointments are scheduled at 1-2 weeks, 6 weeks, 3 months, and 6 months post-surgery.

\begin{itemize}
  \item **Suture/Staple Removal:** Non-dissolvable sutures are typically removed within 7-14 days.
  
  \item **Physical Therapy (PT):** An individualized PT program is initiated shortly after surgery, focusing on pain control, range of motion, and strengthening.

  \item **Home Exercise Program:** Patients are provided with exercises to perform at home to supplement PT.

  \item **Activity Resumption Guidelines:** Clear guidelines are provided for gradually increasing activity levels.

  \item **Imaging:** Routine follow-up imaging is generally not required unless specific concerns arise.

  \item **Warning Signs:** Patients are educated on warning signs that necessitate immediate contact with their surgeon.
\end{itemize}

Adherence to follow-up schedules and rehabilitation plans is essential for achieving optimal recovery and long-term knee health.

\section*{Outcomes \& Prognosis}
Meniscal tears are among the most common knee injuries, affecting a wide demographic, from athletes to the elderly. The estimated incidence ranges from 60 to 70 per 100,000 people annually. Men are more frequently affected, particularly in younger populations engaged in contact sports. Traumatic tears are common in individuals under 40-50 years, while degenerative tears are prevalent in older populations, often coexisting with osteoarthritis. Meniscal surgery, particularly partial meniscectomy, is one of the most frequently performed orthopedic procedures globally. While mortality rates are low, morbidity can include persistent pain, joint stiffness, and the long-term development of osteoarthritis.

Outcomes following meniscal surgery are generally favorable, with variations based on the procedure type and patient factors. For **partial meniscectomy**, success rates for pain relief and functional improvement exceed 80-90\%. However, there is an increased risk of developing knee osteoarthritis due to meniscal tissue loss.

**Meniscal repair** aims to preserve native meniscal tissue, offering a superior long-term prognosis regarding osteoarthritis prevention. Success rates for meniscal repair range from 60-90\%, influenced by factors such as tear location, pattern, patient age, and concomitant injuries. Despite successful repair, some patients may experience residual symptoms. Overall, both procedures aim to improve quality of life by alleviating symptoms, with meniscal repair offering the potential for better long-term knee health.

\section*{References}
\begin{enumerate}
  \item MedlinePlus. Meniscus repair. U.S. National Library of Medicine; 2024. Available from: https://medlineplus.gov/meniscusrepair.html
  
  \item Miller MD, Thompson SR. DeLee and Drez's Orthopaedic Sports Medicine: Principles and Practice. 5th ed. Elsevier; 2020.
  
  \item American Academy of Orthopaedic Surgeons (AAOS). Management of Meniscal Tears. Clinical Practice Guideline; 2024.
  
  \item Canale ST, Beaty JH. Campbell's Operative Orthopaedics. 14th ed. Elsevier; 2021.
\end{enumerate}

\clearpage
